\documentclass[11pt,a4paper]{article}
\usepackage[utf8]{inputenc}
\usepackage[T1]{fontenc}
\usepackage{amsmath,amssymb}
\usepackage{graphicx}
\usepackage{hyperref}
\usepackage[numbers]{natbib}
\usepackage{geometry}
\geometry{margin=1in}

\title{Daily Paper Review: Cortex: Consolidating Context into Parametric Memory}
\author{Auto-generated Report}
\date{October 16, 2026}

\begin{document}
\maketitle

\section{Paper Summary}

\subsection{Problem and Motivation}
Context improves LLM performance, but gains are lost once the context is removed. Zhang et al.\ (arXiv:2610.20399) formalize \textbf{Memory Consolidation (MC)} for LLMs, inspired by biological hippocampal-neocortical systems consolidation: given a model $M$ and a context $C$ that substantially improves performance, MC seeks an updated model $M'$ that (1)~operates \emph{without} $C$, (2)~recovers the context-induced gains on target queries, and (3)~preserves the original model's capabilities outside the intended change. The \textbf{context gap} $\delta_{\text{ctx}}(M, Q_{C^+}, C) := \text{Perf}(M, Q_{C^+}, C) - \text{Perf}(M, Q_{C^+})$ quantifies the behavioral improvement that is lost when $C$ is removed.

This formulation \emph{unifies} previously isolated problems---knowledge editing, skill learning (new APIs/tools), behavioral steering (persona adoption), and knowledge suppression (unlearning)---under a single objective. Existing approaches fail along different dimensions: RAG requires retrieval at inference and misses implicitly relevant knowledge; long context windows provide only ephemeral gains; standard fine-tuning lacks continual-learning safeguards; and specialized editing methods (ROME, MEMIT, AlphaEdit) handle only narrow task types with limited consolidation rates.

\subsection{Method}
Cortex implements MC through three design axes: \emph{what}, \emph{where}, and \emph{how} to consolidate.

\paragraph{What to Consolidate: Prospective Rehearsal.}
Rather than requiring evaluation-time query distributions, Cortex generates \textbf{prospective queries} anticipating future uses of the context:
\begin{equation}
q_i \sim M(\cdot \mid I_q, C), \quad y_i \sim M(\cdot \mid C, q_i), \quad i = 1, \ldots, N
\end{equation}
where $I_q$ is an instruction requesting queries that anticipate uses of $C$. Both queries and responses are generated by the pre-update model itself (self-distillation), reducing distribution mismatch.

\paragraph{Where to Consolidate: FFN Plasticity.}
FFNs are chosen as the memory substrate due to their key-value memory characteristics. Two forms of plasticity operate on bias-free SwiGLU layers $F(x) = W_D[(W_U x) \odot \text{SiLU}(W_G x)]$:

\emph{Synaptic Plasticity:} Low-rank LoRA updates restricted to FFN layers only:
\begin{equation}
W' = W + B_W A_W, \quad W \in \{W_G, W_U, W_D\}
\end{equation}

\emph{Structural Plasticity:} Adds nonlinear feature channels while completely freezing original weights:
\begin{equation}
F'(x) = F(x) + D_s\bigl[(U_s x) \odot \text{SiLU}(G_s x)\bigr]
\end{equation}
where $U_s, G_s \in \mathbb{R}^{k \times d_h}$ and $D_s \in \mathbb{R}^{d_h \times k}$. Both variants fold into standard FFN layers at inference---no separate recall mechanism.

\emph{MoE Extension:} For mixture-of-experts models, a \textbf{shared expansion branch} processes every token rather than per-expert expansions (which suffer from uneven training exposure):
\begin{equation}
F'_{\text{MoE}}(x) = \textstyle\sum_e g_e(x)\bigl(E_e(x) + E_s(x)\bigr)
\end{equation}
The shared branch folds into every expert when routing weights sum to 1.

\paragraph{How to Consolidate: Context-to-Memory Distillation.}
A teacher-student framework where the frozen pre-update model \emph{with} context teaches the student model \emph{without} context. Only new FFN parameters $\psi_t$ are trained:
\begin{equation}
\mathcal{L}_t(\psi_t) = \mathbb{E}_B\!\left[\frac{1}{|J_B|}\sum_{i \in J_B}\left[\beta T^2 \text{KL}_{K_v}(p_i^T \| p_i^S) + \frac{\alpha}{L \cdot d_h}\sum_{l=1}^{L}\|h_{l,i}^S - h_{l,i}^T\|_2^2\right]\right]
\end{equation}
The KL term uses a truncated teacher distribution (top-$K_v$ logits) for storage efficiency. Hidden-state alignment is measured \emph{after the FFN residual update}, motivated by compensating for context removal through FFN corrections.

\paragraph{Continual Consolidation.}
For sequential updates, Cortex retains only earlier \emph{queries} (not contexts or answers), regenerating responses at each update via the current model. A set of generic queries with responses from the initial model $M_0$ prevents capability drift. Earlier adapter parameters are frozen during subsequent updates.

\subsection{Results}
Experiments span 4B--120B parameters across dense (Qwen2.5-14B, Qwen3-4B, Gemma3-12B) and MoE (GPT-OSS-20B, GPT-OSS-120B) architectures, evaluating five task types: factual editing, skill learning (code/API), behavioral steering, knowledge suppression, and agent experience (ALFWorld, $\tau$2-Bench).

\textbf{Single/Batch Consolidation:} On persona steering, Cortex achieves \textbf{97--100\% CR with 99--100\% RR} across models. On novel API learning (Qwen2.5-14B batch), Structural reaches 74\% CR / 98\% RR vs.\ OPCD's 40\% CR / 100\% RR. On ALFWorld agent experience, Cortex achieves \textbf{128--133\% CR} (surpassing contextual prompting) with $\sim$88\% RR.

\textbf{50 Sequential Factual Edits (Qwen2.5-14B):} After 50 edits, Cortex-Structural achieves \textbf{88.9\% CR / 99.8\% RR} with only 0.74M parameters/edit, while ROME reaches 28.5\% CR, MEMIT 23.2\%, AlphaEdit 22.7\%, and WISE 28.0\%. Baselines produce 86--91\% repetition/length-limit termination at checkpoint~50; Cortex produces 0.0--0.8\%. Knowledge retention after 40 subsequent edits: Structural retains 97.2\% of the first 10 facts' CR.

\textbf{Baselines vastly outperformed:} Doc2LoRA achieves 95\% CR on single edits but collapses to 24\% in batch; SHINE has strong single-edit recovery but catastrophic RR (4--10\% in suppression). Context distillation methods (OPCD, SELF-PARAM, SDCD) achieve 19--63\% CR. Test-time training methods (LoCAS, $\delta$-Mem) achieve $\leq$29\% CR.

\subsection{Limitations}
(1)~\textbf{Selective suppression costs}: Suppressing target responses disrupts related non-target responses (Neighbor probes drop 9--20 points). (2)~\textbf{Environment-dependent}: $\tau$2-Bench shows only 38.5\% CR (Structural) on GPT-OSS-20B, indicating consolidation effectiveness varies by task. (3)~\textbf{Batching degrades}: Single-to-batch CR drops can be large (e.g., 91\%$\to$45\% on GPT-OSS-20B code). (4)~\textbf{Portability gap}: Even Structural achieves only 49.4\% portability vs.\ 95.3\% efficacy after 50 edits, showing compositionality remains challenging. (5)~\textbf{Computational cost}: Each consolidation requires forward passes through the full backbone. (6)~\textbf{No structural expansion pruning}: Structural branches fold into wider FFNs but parameter count grows linearly with number of consolidations.

\subsection{Relevance to Research}
Cortex directly addresses Topic T3 (continual learning for LLM agents), providing the mechanistic complement to the activation-focused ExpActivator (Review~\#124): where ExpActivator determines \emph{when} to use accumulated experience, Cortex determines \emph{how} to permanently integrate it into model parameters. The FFN-only plasticity design connects to How LoRA Remembers (Review~\#62), which established parametric memory laws for LoRA fine-tuning---Cortex's finding that FFN-only adaptation outperforms full LoRA (+3.8 CR, +4.0 RR) refines this: \emph{where} you write matters as much as \emph{how much} you write. The prospective rehearsal mechanism resonates with the Escaping the Self-Confirmation Trap paradigm (Review~\#78), which used execute-distill-verify to prevent agents from reinforcing incorrect experiences; Cortex's self-generated queries serve a similar anticipatory function but for parametric consolidation rather than experience replay. The structural plasticity's additive nonlinear channels parallel the approach in Macaron-V1 (Review~\#112), which used mixture-of-LoRA for continual learning; Cortex's structural variant achieves comparable goals with fewer parameters (0.74M/edit) while proving that compact writes outperform larger ones. The query replay for continual consolidation connects to Geometry Conflict (Review~\#49), which analyzed forgetting in LLM continual post-training through gradient geometry---Cortex's finding that freezing earlier parameters does NOT freeze their predictions directly validates the gradient interference mechanism identified in that review.

\section{Follow-up Research Directions}

\subsection{Direction 1: Consolidation-Aware Experience Activation for Hybrid Parametric-Contextual Agents}
\textbf{Gap:} Cortex consolidates \emph{all} context into parametric memory, while ExpActivator (Review~\#124) keeps all experience in an external retrieval store. Neither considers the optimal split: some experiences (stable routines, frequently used knowledge) benefit from parametric consolidation, while others (rare edge cases, recently acquired and not yet verified) are better kept contextual. No existing method dynamically decides which experiences to consolidate and which to keep retrievable.

\textbf{Approach:} Build a two-tier agent memory system. Tier~1 is an ExpActivator-style retrieval store with per-step screen matching. Tier~2 is Cortex-style parametric memory via structural FFN expansion. A \textbf{consolidation trigger} monitors each experience's activation frequency and outcome consistency: when an experience is activated more than $n_{\text{consolidate}}$ times with consistent positive outcomes (measured by ExpActivator's screen-matching score exceeding a threshold), it is promoted from Tier~1 to Tier~2 via Cortex's context-to-memory distillation. The consolidation trigger uses the routine support computation $M_j(x)$ from ExpActivator as a readiness signal: experiences with high temporal regularity and scenario consistency are consolidation-ready. After consolidation, the Tier~1 entry is retained as a fallback (activated only when the parametric prediction confidence drops below a threshold), providing graceful degradation. This creates a biologically-inspired memory lifecycle: experiences begin as episodic (Tier~1) and graduate to semantic (Tier~2) through repeated activation and verification.

\textbf{Feasibility:} High. Both components (ExpActivator and Cortex) are independently validated. The consolidation trigger requires only tracking activation frequency and outcome, which ExpActivator already provides. The main risk is interference: consolidating one experience may degrade parametric memory of previously consolidated experiences, requiring Cortex's query replay mechanism.

\subsection{Direction 2: Memory Consolidation for Diffusion LLM Agents via Denoising-Step-Aware Distillation}
\textbf{Gap:} Cortex's context-to-memory distillation assumes autoregressive generation where teacher and student produce tokens left-to-right. Diffusion LLMs generate text through iterative block denoising with bidirectional attention, fundamentally changing the distillation dynamics: the teacher's hidden states at early denoising steps (high masking) differ qualitatively from late steps (low masking), and the student must learn to compensate for context removal across \emph{all denoising steps simultaneously}. No existing MC method handles diffusion LLM architectures.

\textbf{Approach:} Extend Cortex's distillation loss to be \textbf{denoising-step-aware}. For each training example, run the teacher (with context) and student (without context) through $T$ denoising steps. At each step $t$, compute the KL and MSE losses separately, then weight them by the step's information content: early steps (high masking) contribute more to semantic layout and should receive higher KL weight; late steps (low masking) contribute more to token-level refinement and should receive higher MSE weight. The step-dependent weighting follows the curvature analysis from SpecFold (Review~\#122): positions where the denoising trajectory has high curvature need more careful distillation supervision. For structural plasticity in diffusion LLMs, the added FFN channels should be \emph{step-conditioned}: $F'(x, t) = F(x) + D_s[(U_s x) \odot \text{SiLU}(G_s x)] \cdot w(t)$, where $w(t)$ is a learned step-dependent gating scalar, allowing the consolidated memory to activate differently at different denoising steps.

\textbf{Feasibility:} Medium. The core distillation framework generalizes naturally. The main challenge is computational: diffusion LLMs require $T$ forward passes per denoising trajectory, multiplying distillation cost by $T$. SoloQ's W4A4 quantization (Review~\#123) could mitigate this by running quantized teacher/student forward passes during distillation.

\subsection{Direction 3: Scalable Continual Consolidation via Tracking-Value-Guided Selective Rehearsal}
\textbf{Gap:} Cortex's query replay for continual consolidation replays \emph{all} previously consolidated queries at each new update, with cost growing linearly in the number of past consolidations. The paper tests only 50 sequential edits; scaling to hundreds or thousands---as required for deployment-length agent learning---would make rehearsal prohibitively expensive. Moreover, uniform replay over all past queries ignores that some consolidated memories are more fragile (more likely to be disrupted by future updates) than others.

\textbf{Approach:} Use TPPO's tracking-value concept (Review~\#125) to prioritize rehearsal. For each consolidated query $q_i$, maintain a \textbf{consolidation tracking value} $\hat{V}_i = \sigma(\sum_{j=1}^{t} c_{i,j})$ that estimates the probability that query $q_i$'s consolidated answer will remain correct after the current update. At each new consolidation step, compute the per-query increment $c_{i,t} = f_\phi(h_i^{\text{post-FFN}})$ based on the student's hidden state at $q_i$'s position. Queries with low tracking values (high disruption risk) are prioritized for replay; queries with high tracking values (stable consolidation) are deprioritized or skipped entirely. The replay budget $B_{\text{replay}}$ is fixed regardless of the number of past consolidations, achieving $O(1)$ rehearsal cost per update. The tracking-value critic is itself a lightweight cumulative model (per TPPO's Proposition~4.1) that requires no separate training---it emerges from monitoring the student's hidden-state evolution across consolidation steps.

\textbf{Feasibility:} Medium-high. The tracking-value computation is cheap (one scalar per query per consolidation step). The main risk is that the binary reward assumption (consolidated correctly vs.\ not) may be too coarse---intermediate degradation (partially correct answers) requires the quantile extension proposed in Review~\#125's follow-up Direction~2.

\section{Conclusion}
Cortex establishes Memory Consolidation as a unifying framework for making contextual LLM gains permanent, bridging knowledge editing, skill learning, behavioral steering, and suppression under a single formulation. The dual-plasticity design---synaptic (low-rank FFN updates) and structural (additive nonlinear channels)---achieves 88.9\% consolidation rate with 99.8\% retention after 50 sequential edits using only 0.74M parameters per edit, while baselines collapse to 23--28\% consolidation with 86--91\% degenerate outputs. The most striking finding is that consolidated experience can \emph{surpass} contextual prompting (128--133\% CR on ALFWorld), suggesting that parametric integration enables compositional reasoning that context alone cannot. For continual learning in LLM agents, Cortex provides the missing ``write'' operation complementing retrieval-based ``read'' systems, establishing a deployment-specific learning loop where accumulated experience graduates from contextual to parametric memory.

\begin{thebibliography}{9}
\bibitem{zhang2026cortex}
J.~Zhang, F.~He, X.~Liang, Z.~Xu, H.~Xu, H.~Huang, S.~Wu, and N.~Zhang.
\newblock Cortex: Consolidating Context into Parametric Memory.
\newblock \emph{arXiv preprint arXiv:2610.20399}, 2026.

\bibitem{meng2022rome}
K.~Meng, D.~Bau, A.~Andonian, and Y.~Belinkov.
\newblock Locating and Editing Factual Associations in GPT.
\newblock In \emph{Proc.\ NeurIPS}, 2022.

\bibitem{zhang2026expactivator}
F.~Zhang et~al.
\newblock Relevance Does Not Imply Applicability: Experience Activation for Personal GUI Agents.
\newblock \emph{arXiv preprint arXiv:2609.33304}, 2026.

\bibitem{wang2026tppo}
F.~Wang et~al.
\newblock TPPO: Tracking-Value PPO For Language Model Reasoning.
\newblock \emph{arXiv preprint arXiv:2610.16794}, 2026.

\bibitem{snell2022context}
C.~Snell, D.~Klein, and R.~Rubin.
\newblock Learning by Distilling Context.
\newblock \emph{arXiv preprint arXiv:2209.15189}, 2022.
\end{thebibliography}

\end{document}
